\chapter*{Введение}
\addcontentsline{toc}{chapter}{Введение}
\label{ch:intro}

\textbf{HTML} (HyperText Markup Language) — стандартный язык разметки
веб-документов, описывающий смысловую структуру содержимого: заголовки,
абзацы, списки, таблицы, цитаты, ссылки, изображения. \textbf{CSS}
(Cascading Style Sheets) — язык каскадных таблиц стилей, отвечающий за
визуальное оформление: цвет, типографику, отступы, маркеры списков,
фоновые изображения. Связка <<HTML + CSS>> образует базовый
инструментарий любого современного веб-интерфейса; даже самый сложный
SPA-фреймворк в конечном счёте отрисовывает результат именно через эти
два языка.

В отличие от лабораторных работ~10 и~11, где макеты собирались в
графическом редакторе Figma без единой строчки кода, эта работа
возвращает на уровень разметки и стилей. Вместо локальной среды
используется онлайн-редактор \textbf{Codepen}: каждое из четырёх
заданий — отдельный pen, который студент форкает (`Fork'),
редактирует и сохраняет (`Save'). Преимущество Codepen — мгновенный
preview результата прямо в браузере и возможность сразу поделиться
ссылкой на готовое решение.

Цель работы — освоить базовые элементы языков HTML и CSS на четырёх
независимых практических заданиях, охватывающих разметку текста,
списки и таблицы, селекторы и типографические свойства.

Задачи:
\begin{enumerate}
  \item Освоить семантическую разметку текста: заголовки разных
        уровней, абзацы, цитаты, гиперссылки, изображения с
        атрибутом \texttt{alt} (блок~1.1 методички).
  \item Изучить структуру маркированных и нумерованных списков, а
        также таблиц с группировкой ячеек через \texttt{colspan}
        и \texttt{rowspan} (блок~1.2).
  \item Освоить базовые CSS-селекторы: по тегу, по классу, селектор
        потомка, прямой потомок --- и научиться применять
        свойства \texttt{color}, \texttt{font-style},
        \texttt{font-weight} (блок~2.1).
  \item Освоить типографические свойства CSS для текстовых блоков:
        \texttt{font-size}, \texttt{font-family},
        \texttt{text-decoration}, \texttt{list-style-type},
        \texttt{background-image} (блок~2.2).
\end{enumerate}

\textbf{Стенд:}
\begin{itemize}
  \item Веб-приложение Codepen (\texttt{codepen.io}), доступ через
        современный браузер (Chromium, Firefox); аккаунт создаётся
        через GitHub, Google или электронную почту.
  \item Локальные копии решений в каталоге \texttt{lab12/codepen-solutions/},
        открываются через \texttt{xdg-open} прямо из репозитория без
        авторизации в Codepen.
  \item Внешние таблицы стилей: общий каркас Netology
        \texttt{netology-code.github.io/resourses/codepen/styles.css}
        (для всех четырёх блоков) и дополнительный
        \texttt{style.css} (для блока~2.2); UI-фреймворк
        \texttt{Material~Design~Lite} (\texttt{material.indigo-pink.min.css}).
\end{itemize}

В отличие от лабораторных работ~4--9, стенд не требует ни виртуальных
машин, ни bash-автоматизации. Поэтому в каталоге \texttt{lab12/}
отсутствуют файлы \texttt{config.sh} и сценарии настройки сервисов.

\endinput
